\chapter{Accuracy and Efficiency}
\label{chap:10}
In theory, the algorithms presented in this book are all exact. However, the presence of round-offerrors in the calculations ensures that the results are rarely so. Other sources of error include the use of numerical integration, and the differences in behaviour between a rigid-body system and the physical system it represents. Section 10.1 examines these various sources of error, paying particular attention to the sources of round-offerror in dynamics calculations and ways to avoid them. If the acceleration of a rigid-body system changes greatly in response to a tiny change in the applied forces, then it is said to be sensitive. If you want to know the acceleration of such a system accurately, then the forces must be known even more accurately. It turns out that many kinematic trees have this property, and the phenomenon is examined in Section 10.2. Section 10.3 explores the topic of efficiency. It explains how the efficiency of a dynamics algorithm is measured; and it presents tables and graphs comparing the efficiencies of several versions of the three main algorithms for kinematic trees. The efficiency of the composite-rigid-body algorithm is examined in detail, partly to show how an operations-count formula is obtained, and partly because the efficiency of this particular algorithm varies greatly with the amount of branching in the tree. This section also shows that the true asymptotic complexity of the composite-rigid-body algorithm is O(nd), where d is the depth of the connectivity tree, and that the true complexity of the so-called O(n3) algorithm is O(nd2). One of the disadvantages of model-based algorithms is that they have to cater for the general case. As such, they perform many calculations that are indeed necessary in the general case, but are frequently not necessary for some particular system model. One solution to this problem is the automatic generation of customized dynamics routines—stripped-down versions of the general algorithm that perform only those calculations that are necessary for a specific given system model. Section 10.4 presents a brief description of a technique for doing this, which is called symbolic simplification.

\section{ources of Error Dynamics calculations are subject to three kinds of error: round-offerror, which is the result of using finite-precision floating-point arithmetic; truncation error, which is the result of approximating an infinite series by a finite number of terms; and modelling error, which refers to the discrepancy between the behaviour of the mathematical model and the behaviour of the physical system it represents. Each kind of error is distinct from the others, so let us examine them separately. Round-offError Round-offerrors occur during most floating-point operations, but they are usually tiny. Large errors can arise if two large-magnitude numbers are added or subtracted, producing a small-magnitude result, and especially if this happens several times during the course of a calculation. The simplest defence against round-offerror is to use double-precision floating-point arithmetic for all calculations, but this will not solve every problem. Large round-offerrors can arise during dynamics calculations for a variety of reasons. Some of the more likely ones are: 1. using a far-away coordinate system; 2. large velocities; 3. large inertia ratios; and 4. nearly singular loop-closure constraints. A coordinate system is far away from a rigid body if the distance between the origin and the centre of mass is many times the body’s radius of gyration. Ideally, this distance should be no more than about one or two radii of gyration. This effect is the reason why dynamics algorithms are more accurate when they are implemented in body coordinates rather than absolute coordinates. The large-velocity problem arises in systems with large mean velocities, like spacecraft, in which the relative velocities of the component parts are small compared with their absolute velocities. This problem can be solved by using an inertial reference frame with a velocity close to the mean linear velocity of the system. Large differences in inertia arise from relatively small differences in scale. For example, if two solid spheres are made of the same substance, and one is 10 times larger than the other, then the larger sphere will have 1000 times the mass of the smaller one, and 100,000 times the rotational inertia. Let us examine in more detail how some of these round-offerrors arise. Figure 10.1 shows a planar rigid body having a centre of mass at the point C, which is located relative to the origin by the vector c = −−→ OC = [cx cy]T. We shall use planar vectors in this example, to keep the matrices small, but the round-off problems are essentially the same for planar and spatial vectors. This body has unit mass, and it has unit rotational inertia about C, so its radius of gyration}
\label{sec:10.1}
O x y C ω=1 m = 1 IC = 1 fC c = cx cy fC = 1 −1 1 = τ fCx fCy vC = 1 0 0 = ω vCx vCy Figure 10.1: Calculating the acceleration of a planar rigid body is 1. Expressed in C coordinates, the velocity of this body is vC = [1 0 0]T, and it is subject to an applied force of fC = [1 −1 1]T. The former is a unit rotation about C, and the latter is a pure force having a magnitude of √ 2 and a line of action passing through the point (0.5, 0.5) relative to C. Expressed in O coordinates, the equation of motion for this body is IO aO = fO −vO ×∗IO vO , (10.1) where vO =

1 cy −cx

, vO×∗=

0 cx cy 0 0 −1 0 1 0

, fO =

1 + cx + cy −1 1

and IO =

1 + c2 x + c2 y −cy cx −cy 1 0 cx 0 1

. The exact solution to this equation is aO =

1 cy −1 1 −cx

. Suppose we set cx = 100 and cy = 10, and ask the computer to evaluate aO using IEEE double-precision floating-point arithmetic, which is accurate to one part in 1016. The computer will do this by solving the linear equation

10101 −10 100 −10 1 0 100 0 1

aO =

111 −1 1

, and the calculated value will differ from the exact value by about one part in 1012. Thus, four significant figures of accuracy have been lost. The reason

10 −1 10 0 10 1 10 2 10 3 10 −16 10 −14 10 −12 10 −10 Figure 10.2: Round-offerror in the value of aO, measured at C, versus |c| for this loss is that the coefficient matrix is highly ill-conditioned, having a condition number of 1.02 × 108. To give an idea of how the round-offerror varies with |c|, Figure 10.2 plots ∥aC −CXOaO∥∞against |c|, where aC is the exact value and aO the computed value. It can be seen that the round-offerror is approximately constant when |c| < 1 (i.e., when |c| is less than the radius of gyration), and grows approximately in proportion to |c|2 when |c| > 1. In this example, the linear velocity of the body’s centre of mass was zero (vCx = vCy = 0). If instead the body were to have a large linear velocity, then there could be substantial round-offerror in calculating the right-hand side of Eq. 10.1, even if |c| is relatively small. Round-offError in Dynamics Algorithms Round-offerror is known to increase with body count, and to increase more quickly with some algorithms than with others. The three main algorithms for kinematic trees are known to remain accurate, even when there are thousands of bodies in the system; and it is also known that the articulated-body algorithm is more accurate than the composite-rigid-body algorithm. However, the roundoffbehaviour of many other algorithms is not known. Two relevant articles on this subject are Ascher et al. (1997) and Featherstone (1999b). Truncation Error Truncation errors are the result of using finite approximations to infinite series, continuous curves, continuous functions and the like. In rigid-body dynamics, the main source of truncation error is the numerical integration process, which is only an approximation to exact mathematical integration. By its nature, the management of truncation error involves a compromise between efficiency

and accuracy: a better approximation will reduce the error, but will be more expensive to calculate. The management of integration error involves making an appropriate choice of integration method and time step. If one plots the acceleration variables of a simulated rigid-body system against time, then it will often be the case that the resulting graph contains large spikes: brief periods when the accelerations are much higher than at other times. Under these circumstances, it may be beneficial to use an integration method that automatically adapts its step size to maintain a prescribed accuracy. Such methods take big steps when not much is happening, and tiny steps when the accelerations are high. This strategy can lead to an improvement in the overall efficiency of a simulation run by reducing the total number of times that the dynamics has to be calculated. Modelling Error Modelling errors are those that cause the behaviour of the mathematical model to differ from the behaviour of the physical system. If the mathematical model is a rigid-body system, then the main sources of modelling error are: • the rigid-body assumption, • the ideal-joint assumption (exact kinematic constraints), • unmodelled dynamics, and • inaccuracies in model parameters. In reality, there is no such thing as a truly rigid body, and real joints suffer from imperfections such as friction, compliance and play. Unmodelled dynamics refers to aspects of the physical system that have been ignored in formulating the model. For example, one might choose to ignore the dynamics of the chain in a chain-and-sprocket drive, or one might be forced, by lack of data, to ignore temperature-related or position-dependent variations in the friction at an important joint bearing. Most model parameters refer either to geometry or to inertia. The geometric parameters of a mechanical system are usually known to a reasonable accuracy, but the inertia parameters may be less accurate, or even unknown. If the system can be dismantled, then one can measure the inertias of the individual parts; otherwise, one can use parameter-identification techniques, such as those described in Khalil and Dombre (2002). 10.2 The Sensitivity Problem Some rigid-body systems can be extremely sensitive to small changes in applied forces, or small changes in model parameters. Sensitivity to geometric parameters is only to be expected in closed-loop systems having special geometries, or general geometries that are close to being special. However, the problem is

75 x y Figure 10.3: A 6-link planar mechanism in a zigzag configuration broader than this, and, in particular, it is a problem that affects kinematic trees as well as closed-loop systems. To illustrate the problem, consider the 6-link planar robot mechanism shown in Figure 10.3. Each link has unit length, unit mass, a centre of mass half way along its length, and a rotational inertia about its centre of mass of 1/12. The joints are revolute, and the joint angles alternate between +75◦and −75◦, as shown. The robot is initially at rest, and there are no forces acting on it other than the joint forces. The equation of motion is therefore τ = H¨q . Suppose we wish to impart a desired acceleration of ¨qd = [1 1 1 1 1 1]T to this robot. The exact force required to do this is τd = H¨qd =

126.4936 · · · 97.4663 · · · 69.9762 · · · 43.7998 · · · 21.9371 · · · 6.1646 · · ·

. If we apply exactly this force to the robot, then we will get exactly the desired acceleration. Now suppose that the robot’s actuators are slightly imperfect, such that the forces they actually apply are equal to the commanded forces rounded to three significant figures. The force that will actually be applied to the robot is then τa =

126 97.5 70.0 43.8 21.9 6.16

,

and the robot’s response to this force will be ¨qa = H−1τa =

0.6952 1.3654 1.3808 0.5894 0.9057 1.0705

. Observe that some of the accelerations are as much as 40\% offtheir correct values, despite the applied forces all being accurate to better than 0.5\%. In effect, the mechanism has amplified the error by more than a factor of 80. This problem has been studied in Featherstone (2004), where the condition number of H is used as an indication of the underlying sensitivity of the mechanism. This paper shows that the condition number can grow in proportion to the fourth power of the number of bodies, and that the worst-case condition number for a chain of identical links, like the mechanism in Figure 10.3, is approximately 4N 4 B. The worst-case condition number for the 6-link robot is therefore approximately 5184, although the actual condition number of H in the above example is 725; and the worst case for a 10-link robot is 40,000. Numbers like this indicate that one should be cautious of simulating long chains of rigid bodies, and skeptical of the accuracy of the results. This paper also investigates branched kinematic trees, and it shows that sensitivity increases with the depth of the tree, other things being kept equal. Thus, a branched kinematic tree will usually have a lower condition number than an unbranched tree made from the same set of bodies. The obvious implication is that one should choose a minimum-depth spanning tree for a closed-loop mechanism. 10.3 Efficiency To compare the efficiencies of different algorithms, it is necessary to have a common measure of efficiency. In the case of dynamics algorithms, the most widely used measure is the floating-point operations count. This count is usually broken down into two subtotals: the number of multiplications and divisions, which are collectively referred to as multiplications, and the number of additions and subtractions, which are collectively referred to as additions. Furthermore, the subtotals are expressed as functions of a parameter that measures some aspect of the size of the rigid-body system, such as the number of bodies or the number of joint variables. Some examples of operation-count formulae are shown in Table 10.1. The symbols m and a stand for multiplication and addition, respectively, and n is the number of joint variables. To obtain figures like this, it is necessary to agree on a standard set of rigid-body systems, one for each value of n, so that every algorithm is applied to the same set of systems. The most widely used standard

Algorithm Cost Source RNEA (93n −108)m + (81n −100)a 3 RNEA (93n −69)m + (81n −66)a 1, 3 RNEA (130n −68)m + (101n −56)a 6 RNEA (150n −48)m + (131n −48)a 8 CRBA (10n2 + 22n −32)m + (6n2 + 37n −43)a 7, 14 CRBA (10n2 + 31n −41)m + (6n2 + 40n −46)a 6 CRBA (10n2 + 41n −51)m + (6n2 + 52n −58)a 9, 10 CRBA (11n2 + 11n −42)m + ( 11 2 n2 + 65 2 n −54)a 2, 4 CRBA (12n2 + 56n −27)m + (7n2 + 67n −53)a 12 F\&S ( 1 6n3 + 3 2n2 −2 3n)m + ( 1 6n3 + n2 −7 6n)a 13 ABA (224n −259)m + (205n −248)a 11 ABA (250n −222)m + (220n −198)a 5 ABA (300n −267)m + (279n −259)a 6 Abbreviations RNEA recursive Newton-Euler algorithm CRBA composite-rigid-body algorithm F\&S factorize and solve equation of motion H ¨q = τ −C ABA articulated-body algorithm Sources 1 Balafoutis et al. (1988), Algorithm 3 in Table II 2 Balafoutis and Patel (1989), Algorithm I in Table II 3 Balafoutis and Patel (1991), Algorithm 5.7 in Tables 5.1 and 5.2 4 Balafoutis and Patel (1991), Algorithm 6.2 in Table 6.2 5 Brandl et al. (1988) 6 Featherstone (1987), Table 8-6 7 Featherstone (2005), Eq. 20 with D0a = n −1, D1a = (n2 −n)/2 and D0b = D1b = 0 8 Hollerbach (1980), Newton-Euler method in Table I 9 Lilly and Orin (1991), method IV in Table 7 10 Lilly (1993), method IV in Table 3.7 11 McMillan et al. (1995), Table II (fixed base) 12 Walker and Orin (1982), method 3 13 Table 6.6 with Eq. 6.28 14 Eq. 10.3 Table 10.1: Operation counts for several published dynamics algorithms

2 4 6 8 10 12 14 16 18 0 2000 4000 6000 8000 10000 12000 Comparative costs of RNEA, ABA and O(n3) algorithm RNEA ABA O(n3) algorithm: CRBA+RNEA+F\&S 5 10 15 20 25 30 2 0 5000 10000 15000 Component parts of O(n3) algorithm RNEA F\&S CRBA Total Figure 10.4: Operation counts versus n for various dynamics algorithms is the set of unbranched kinematic chains in which all the joints are revolute, and all the bodies have general inertia and geometric parameters. For easy reference, let us call this set GU (general unbranched), and call its members GU(1), GU(2), and so on. GU(n) is then the unbranched chain with n joint variables, and therefore also n joints and n moving bodies (i.e., NJ = NB = n). Given that the joints in GU are revolute, it is understood that one sine and one cosine calculation are required for each joint. As these calculations are the same for all algorithms that need them, they are omitted from the operation counts. The algorithms that need them are the recursive Newton-Euler and articulated-body algorithms. To give a visual impression of the relative efficiencies of the algorithms, Figure 10.4 shows the total operation count for the best version of each algorithm, plotted against n. It also plots the three component parts of the O(n3) algorithm, so that their relative contributions can be gauged. Strictly speaking, the curves in these graphs are only correct if the cost of performing an addition is the same as the cost of performing a multiplication. However, in practice

it makes little difference whether one costs an addition as one multiplication, 0.5 multiplications or 0.25 multiplications: apart from the vertical scales, the graphs look almost the same. According to these graphs, the O(n3) algorithm is slightly faster than the articulated-body algorithm for n ≤8, and has risen to about 1.6 times the cost of the articulated-body algorithm by the time n = 18. The second graph shows that the cost of the composite-rigid-body algorithm (which is O(n2)) is the dominant term in the O(n3) algorithm from n = 8 until well past n = 30. (The F\&S curve crosses the CRBA curve at approximately n = 45.) It also shows that the cost of factorizing and solving the equation of motion is actually the smallest term when n ≤18. The contribution of the n3 component to the total cost is therefore relatively insignificant at small values of n. 10.3.1 Optimization The figures quoted in Table 10.1 are the result of extensive optimizations, which are aimed at showing each algorithm in its best possible light. The most important of these optimizations are as follows.1 1. Use Denavit-Hartenberg coordinate frames. (See Figure 4.8.) This enables optimizations 2 and 3. 2. Exploit the fact that Si = [0 0 1 0 0 0]T in link coordinates. 3. Implement link-to-link coordinate transformations as a sequence of two axial-screw transforms. An axial-screw transform is a combination of a rotation about and translation along a single coordinate axis. Such transforms are particularly cheap to apply, and formulae and costings are presented in Section A.4. A transformation from one Denavit-Hartenberg coordinate frame to another can be accomplished by two axial-screw transforms: one about the x axis, and one about the z axis. The former is a constant, while the latter includes the joint variable. To see how these optimizations are applied, let us calculate the operationcount formula for the composite-rigid-body algorithm. Table 10.2 reproduces the pseudocode of this algorithm (from Table 6.2), together with some annotations. The symbols ra, rx and vx represent the cost of adding two rigid-body inertias, transforming a rigid-body inertia and transforming a vector, respectively. We shall work out these costs in a moment. The first step is to identify which lines contain floating-point calculations. Assignment and transpose operations do not count as calculations, so there are no floating-point calculations on Lines 1, 3 or 16. Using optimization 2, the 1One more optimization that ought to be mentioned, although it does not apply to the set GU, is to use planar vectors for calculating the dynamics of planar rigid-body systems. This is particularly beneficial for the articulated-body algorithm, as the number of independent parameters in an articulated-body inertia shrinks from 21 to 6.

Algorithm: 1 H = 0 2 for i = 1 to NB do 3 Ic i = Ii 4 end 5 for i = NB to 1 do 6 if λ(i)̸ = 0 then 7 Ic λ(i) = Ic λ(i) + λ(i)X∗ i Ic i iXλ(i) 8 end 9 F = Ic i Si 10 Hii = ST i F 11 j = i 12 while λ(j)̸ = 0 do 13 F = λ(j)X∗ j F 14 j = λ(j) 15 Hij = F TSj 16 Hji = HT ij 17 end 18 end Annotations: loop 1 loop 2 operations: ra + rx F = column 3 of Ic i Hii = element 3 of F loop 3 operations: vx Hij = element 3 of F Table 10.2: Calculations performed by the composite-rigid-body algorithm multiplication Ic i Si on line 9 simplifies to extracting column 3 from Ic i , which does not require any calculation. Likewise, the multiplications ST i F and F TSj on lines 10 and 15 both simplify to extracting element 3 from F , which does not require any calculation. Thus, there are only two lines where floating-point calculations do occur: line 7, which performs the operations ra and rx, and line 13, which performs the operation vx. If this algorithm is applied to the mechanism GU(n), then the body of loop 2 will be executed n times. The condition λ(i)̸ = 0 will be true for every value of i except i = 1, so line 7 will be executed a total of n −1 times. On each iteration of loop 2, the body of loop 3 will be executed i −1 times (because λ(j) = j −1); so line 13 will be executed a total of Pn i=1(i −1) times, which works out as n(n −1)/2 times. The cost formula for the composite-rigid-body algorithm is therefore cost = (n −1)(ra + rx) + n(n −1) 2 vx . (10.2) The next step is to express ra, rx and vx in terms of m and a. If rigid-body inertias are stored in compact form, as explained in Section A.3, then the cost of summing two inertias is ten floating-point additions; so ra = 10a. To obtain the lowest operation counts for rx and vx, we use optimizations 1 and 3. From the figures in Section A.4, the cost of a single axial-screw transform of a spatial

vector is 10m+6a, but vx requires two, so vx = 20m+12a. The basic cost of an axial-screw transform of a rigid-body inertia is 16m+15a, to which 2m+3a must be added each time the angle changes. However, we can shave 1m offthis figure by precomputing the products of the masses with the distance parameters of the screw transforms, since they are all constants. This results in a cost figure of 2(15m + 15a) + 2m + 3a = 32m + 33a for rx. Substituting these costs into Eq. 10.2 gives cost = (10n2 + 22n −32)m + (6n2 + 37n −43)a . (10.3) This formula is as far as we will take the optimization process, but there are further improvements in efficiency that could be made. For example, modest reductions in the costs of both ra and rx can be achieved by modifying the inertia parameters as described in Section 9.7. Another cost saving can be obtained by observing that only element 3 of F is needed to calculate Hij, so the last execution of line 13 in each invocation of loop 3 could be replaced by a simpler calculation that only calculates element 3 of the transformed vector. Even this is not yet the last word in efficiency. However, it is not really worthwhile to implement such fiddly optimizations manually. A better strategy is to get the computer to do the optimization for you, and that is the subject of Section 10.4. 10.3.2 Branches One of the problems with using GU as the benchmark set of rigid-body systems is that it does not contain any branched trees, and therefore does not allow the effect of branches to be seen. It turns out that branches have very little effect on the recursive Newton-Euler and articulated-body algorithms, but they can greatly reduce the cost of the composite-rigid-body algorithm, as well as the cost of factorizing and solving the equation of motion. In both cases, the cost savings are caused by branch-induced sparsity in the joint-space inertia matrix. (See §6.4.) It therefore follows that the cost formulae in Table 10.1 and the curves in Figure 10.4 present the O(n3) algorithm in an unduly pessimistic light, by stating only its worst-case performance. Indeed, even the name ‘O(n3) algorithm’ is misleading, since the true complexity of this algorithm is O(nd2), where d is the depth of the connectivity tree. If there are no branches, then d = n; but if there are branches, then d can be much smaller than n; and if d is subject to a fixed upper limit, then the complexity of this algorithm is only O(n). Let us recalculate the cost formula for the composite-rigid-body algorithm, assuming that it is being applied to a general kinematic tree with n revolute joints, and that the joint axes satisfy Si = [0 0 1 0 0 0]T in link coordinates. Referring back to Table 10.2, it is still the case that lines 7 and 13 are the only two that perform floating-point calculations, but the number of times these lines get executed has changed. In the case of line 7, the expression λ(i)̸ = 0

is false for every body i that is a child of body 0; that is, for every i satisfying i ∈µ(0). Line 7 is therefore executed n −|µ(0)| times. In the case of line 13, loop 3 iterates over every ancestor of body i, starting with body i itself, and terminating on the ancestor nearest the root (which is the ancestor satisfying j ∈µ(0)). The body of loop 3 is therefore executed |κ(i)| −1 times on each iteration of loop 2; and so line 13 is executed a total of Pn i=1(|κ(i)| −1) times. Let us define the quantities D0 and D1 as follows: D0 = n −|µ(0)| (10.4) and D1 = n X i=1 (|κ(i)| −1) . (10.5) The cost of the composite-rigid-body algorithm can then be expressed as cost = D0(ra + rx) + D1vx . (10.6) Furthermore, as |µ(0)| is bounded by 1 ≤|µ(0)| ≤n, and |κ(i)| by 1 ≤|κ(i)| ≤i, it follows that D0 and D1 are bounded by 0 ≤D0 ≤n −1 (10.7) and 0 ≤D1 ≤n(n −1)/2 . (10.8) Both quantities reach their lower limit in a system where every body is connected directly to the base. In such a system, H is diagonal and constant, and the cost of calculating it is zero. At the other extreme, both quantities reach their upper limit in an unbranched kinematic tree. In this case, Eq. 10.6 is identical to Eq. 10.2. If d is the depth of the connectivity tree, then |κ(i)| is also bounded by |κ(i)| ≤d, which implies 0 ≤D1 ≤n(d −1) . (10.9) Thus, the true complexity of the composite-rigid-body algorithm is actually O(nd), and the expression O(n2) only applies in the worst case. If we combine this with the O(nd2) complexity of the factorization algorithms in Section 6.5, then the complete forward-dynamics algorithm has a complexity of O(nd2). Let us now express the cost formula in Eq. 10.6 in terms of m and a. At this point, a new complication arises: only one child per nonterminal body can benefit from the cost savings afforded by Denavit-Hartenberg coordinate frames. We therefore rewrite Eq. 10.6 as follows: cost = D0ra + Ddh 0 rxdh + Dg 0rxg + Ddh 1 vxdh + Dg 1vxg , (10.10) where the superscripts dh and g refer to Denavit-Hartenberg and general coordinate transforms, respectively. So rxdh is the cost of calculating λ(i)X∗ i Ic i iXλ(i)

when iXλ(i) is a Denavit-Hartenberg transform (two axial screws); vxg is the cost of calculating λ(j)X∗ j F when λ(j)X∗ j is a general coordinate transform; Dg 0 is the number of times that rxg is performed; and so on. Note that Ddh 0 +Dg 0 = D0 and Ddh 1 + Dg 1 = D1. To obtain expressions for the various D quantities, we introduce two new sets: πdh and πg. πdh is the set of all bodies i such that λ(i)̸ = 0 and iXλ(i) is a Denavit-Hartenberg transform; and πg is the set of all bodies i such that λ(i)̸ = 0 and iXλ(i) is a general transform. Thus, πdh ∪πg is the set of all bodies that are not children of the root node (i.e., not elements of µ(0)), and |πdh ∪πg| = n −|µ(0)| = D0. Given these sets, we can define the D quantities as follows: Ddh 0 = |πdh| , Ddh 1 = X i∈πdh |ν(i)| , Dg 0 = |πg| , Dg 1 = X i∈πg |ν(i)| . (10.11) The expressions for Ddh 1 and Dg 1 are obtained as follows. When we first evaluated D1, we asked the question ‘How many times is loop 3 executed?’ Another approach would be to ask how many times λ(j)X∗ j is used, for each possible value of j. The answer to this question is that if λ(j) = 0 then this transform is never used; otherwise, it is used once for each body in ν(j), which means it is used |ν(j)| times. The expressions in Eq. 10.11 are simply the sums of these usage counts. The final step is to express ra, rxdh, etc., in terms of m and a. From Sections A.3 and A.4, these expressions are: ra = 10a, rxdh = 32m+33a, rxg = 47m+48a (three axial screws), vxdh = 20m+12a and vxg = 24m+18a. Substituting these expressions into Eq. 10.10 gives cost = Ddh 0 (32m + 43a) + Dg 0(47m + 58a) + Ddh 1 (20m + 12a) + Dg 1(24m + 18a) . (10.12) This is the cost of the composite-rigid-body algorithm when it is applied to a general kinematic tree in which all the joints are revolute. On an unbranched tree, we have Ddh 0 = n −1, Ddh 1 = n(n −1)/2 and Dg 0 = Dg 1 = 0, which makes Eq. 10.12 the same as Eq. 10.3. More on this subject, including a cost formula for floating-base kinematic trees, can be found in Featherstone (2005). To give an idea of how big a difference branches can make, Figure 10.5 reproduces the cost curves from Figure 10.4 for the articulated-body and O(n3) algorithms, the latter now labelled O(nd2), and it adds one new curve: the cost of applying the O(nd2) algorithm to a binary tree. This curve is almost identical to the curve for the articulated-body algorithm, and very different from the curve for the O(nd2) algorithm applied to an unbranched chain. The reason for this is that d is O(log(n)) on a binary tree, so the complexity of the O(nd2) algorithm is only O(n log(n)2) on the binary tree, which is very close to being O(n).

2 5 10 15 20 25 30 35 40 45 50 0 0.5 1 1.5 2 x 10 4 Effect of branches on cost of O(nd2) algorithm ABA on unbranched chain O(nd2) on unbranched chain O(nd2) on binary tree Figure 10.5: Effect of branches on the cost of the O(nd2) algorithm 10.4 Symbolic Simplification The advantage of model-based dynamics software is that each algorithm can be implemented once, and the resulting code will work on any rigid-body system that is supplied to it as an argument. The disadvantage is that the code must cater for every possibility that is allowed by the system model, and cannot take short cuts that depend on prior knowledge of the system. In particular, many of the parameters in a typical system will be zero, and the model-based software will spend much time calculating terms that are bound to evaluate to zero because one of their factors is a zero-valued model parameter. The alternative is to write custom code that is designed to work only on a single rigid-body system, or a narrow class of similar rigid-body systems. Such code can exploit prior knowledge of joint types, geometric and inertia parameter values, and so on, with the end result that fewer arithmetic operations are required to perform the calculation, as compared with a model-based coding of the same algorithm. The manual construction of custom code is not practical, except for the simplest of rigid-body systems; but the automatic construction of such code is entirely practical, and the procedure for doing so is called symbolic simplification. Figure 10.6 illustrates the process. In the first stage, an algorithm and a system model are fed to a symbolic executer. This is a program that applies the algorithm to the data, but, instead of actually performing the numerical calculations, it merely records them, and outputs the record as a data file. The contents of this file are a list of assignment statements detailing the floatingpoint calculations that the given algorithm would have performed on the given model if it had been executed for real. This file also identifies the input and output variables. In the second stage, the assignment-statement file is passed to a symbolic simplifier program. This program reads the file, applies a few basic simplifica-

system model algorithm symbolic executer symbolic simplifier assignment statement list custom computer code Figure 10.6: The symbolic simplification process tion rules, and outputs the result in the form of computer source code expressed in a suitable programming language. The most important simplifications are as follows: • Remove any assignment statement that calculates a quantity that is never used. (Being output counts as being used.) • Remove assignment statements of the form a = b, where a is not an output variable, and b is either a constant or a variable, and replace all subsequent instances of a with b. • Perform some basic algebraic simplifications, such as: x·0 →0, x·1 →x, x + 0 →x, and so on. • Evaluate constant expressions. These rules must be applied exhaustively, as the application of one rule can create new opportunities for the application of another. They must also be cheap, as there can easily be tens of thousands of assignments in the list. Table 10.3 presents an example of how the process works. In this example, v is an input variable, τout is an output variable, and I and s are quantities defined in the system model. The algorithm specifies how to calculate τout from v. If we were to implement this algorithm directly, then it would require one 3 × 3 × 1 matrix multiplication, one 3D vector cross product and one 3D vector dot product. This adds up to a total of 29 floating-point arithmetic operations, which can be seen in the assignment-statement file. The symbolic executer deals only with scalars, so it expands vector operations into their scalar equivalents. It also processes only one vector operation at a time, so it introduces temporary variables as required. Apart from that, the only other thing it does is to substitute model parameters with their actual values. Using only the simplifications listed above, the simplifier is able to cut the total operations count from 29 down to 5. If the simplifier had used a more

calculation specified in the algorithm p = v × I v τout = p · s data supplied in the model I =

1.5 0 0 0 1.6 0 0 0 1.7

, s =

0 0 1

contents of the assignment-statement file tmp001 = 1.5 * v(1) + 0 * v(2) + 0 * v(3) tmp002 = 0 * v(1) + 1.6 * v(2) + 0 * v(3) tmp003 = 0 * v(1) + 0 * v(2) + 1.7 * v(3) p(1) = v(2) * tmp003 - v(3) * tmp002 p(2) = v(3) * tmp001 - v(1) * tmp003 p(3) = v(1) * tmp002 - v(2) * tmp001 tau\_out = p(1) * 0 + p(2) * 0 + p(3) * 1 contents of the source-code file tmp001 = 1.5 * v(1) tmp002 = 1.6 * v(2) p(3) = v(1) * tmp002 - v(2) * tmp001 tau\_out = p(3) Table 10.3: Symbolic simplification example complex and sophisticated set of rules, then it might have found the optimal solution, which is τout = 0.1v1v2, but the process would have required much more CPU time. Even so, the simple rule set is enough to find 90\% of the possible cost savings. In this example, the simplifier cut the cost by almost a factor of 6. In practice, this factor is highly sensitive to the number of zeros in the system model. Factors as high as 10 have been reported in the literature (Wittenburg and Wolz, 1985); but factors close to 1 are also possible. Several variations are possible on the above scheme. For example, a set of predefined algorithms can be built into the symbolic executer, so that the user only has to supply the system model. This makes the software easier to use, and easier to implement, but it removes the user’s freedom to supply an algorithm of his own. Another variation is that the symbolic executer and simplifier could be two parts of the same program, rather than two separate programs. In this case, the assignment-statement file would not exist, although an equivalent internal data structure would still be required. Another possibility is to allow the model to define only some of the parameter values. The undefined parameters would then be treated as unknown quantities to be supplied at run time. This feature allows the creation of custom code that caters for a set of rigid-body systems, each differing from the others only in the values of specific parameters. One potential disadvantage of symbolic simplification is that it reduces vector operations to scalar operations. On a computer with vector-arithmetic hardware, the resulting code may make poor use of the available processing power,

and could end up slower than the original vector-oriented code. Another possible disadvantage is that the size of the customized code grows linearly with the operations count (because it is just a list of assignment statements). For a sufficiently complicated rigid-body system, the custom code will be too large to fit in the computer’s instruction cache. When this happens, the computer is forced to fetch instructions from main memory, which is a much slower process. The net effect is that execution speed goes down because the CPU is repeatedly having to wait for the next batch of instructions. Several programs have been written that implement this technique, or something similar, and two have become commercially available: SD/FAST and Autolev.2 More on the subject of automatic generation and simplification of dynamics equations can be found in Murray and Neuman (1984); Neuman and Murray (1987); Wittenburg and Wolz (1985). 2At the time of writing, further information on these two programs could be found at http://www.autolev.com and http://www.sdfast.com.
